\subsection{Selected quadratic companions and a faithful free factor}
\label{sec:selected-pair-faithfulness}

We retain the arithmetic cyclotomic kernel $H$, its actual rational
outer action $\rho$, and the notation of
Theorem~\ref{thm:odd-torsion-cubic}. In particular
$M=H^{\mathrm{ab}}$, $T=M_{\Lambda\text{-tors}}$,
$Q=M/T$, and $R=\Lambda[1/p]$. The unconditional structural
calculations give $H$ free pro-$p$, $T$ finite free over
$\mathbf Z_p$, and $Q[1/p]$ free of rank one over $R$ in each
odd $\Delta$-character, with no even component. The rational
target is the weight-completed free prounipotent group on the
odd generators, as established earlier. No assertion about all
the remaining torsion companions is imposed in the next result.

\begin{theorem}[K17-SELECTED-PAIR]
\label{thm:selected-pair-faithfulness}
Let $p\geq5$, and let $I$ be a set of even residue classes
modulo $p-1$ such that
\[
                 X^{[e]}=X^{[2-e]}=0\qquad(e\in I).
 \tag{SP.1}
\]
Then, for every $e\in I$, the even module $T_e=T^{[e]}$
is rationally cyclic over $R$ and its actual free-$Q$ quadratic
companion satisfies Pair-NV at every root, with repeated roots
allowed. For $e=0,2$ the module $T_e$ is zero, so this condition
is empty.

Choose a $\Delta$-equivariant continuous $\mathbf Z_p$-splitting
of $M\to Q$, a convergent product basis of $Q$, and finite
bases of the character components of $T$. Choose compatible
exact eigenlifts forming an actual compact free basis of $H$,
as in the earlier basis construction. Let $H_I$ be the closed
subgroup generated by all the chosen $Q$ lifts and the chosen
basis lifts of $\bigoplus_{e\in I}T_e$. Then $H_I$ is a closed
free factor of $H$, and
\[
                         \rho|_{H_I}\text{ is injective}.
 \tag{SP.2}
\]
No cyclicity or cubic nonvanishing condition is required on
the omitted components. The subgroup is specified by these
basis choices; it need not be $\Gamma$-stable.
\end{theorem}

\begin{proof}
\par\medskip\noindent\textbf{The two distinct character inputs.}
Fix $e\ne0,2$ in $I$ and put
\[
              \chi=\omega^{1-e},\qquad
              \theta=\omega^{2-e}.
 \tag{SP.3}
\]
First apply Wake--Wang-Erickson's Proposition~2.2.6 to the
reflected even character $\omega^e$. Its corresponding odd
character is $\omega^{e-1}$, whose inverse is $\chi$.
The first zero in (SP.1) therefore makes $M_\chi$ integrally
free of rank one over $\Lambda$. The natural unramified
quotient $M_\chi\twoheadrightarrow X_\chi$ shows that
$X_\chi$ is cyclic. The involution and Tate twist preserve
cyclicity, so the Kummer pseudoisomorphism
\[
                         M^+\sim(X^-)^\iota(1)
\]
makes $T_e[1/p]$ cyclic over $R$. Finite kernels and cokernels
disappear upon inverting $p$. The exact component is $e=1-(1-e)$,
including the Tate twist; see
\cite[Proposition~3.4.7 and Corollary~3.4.8,
p.~92]{SharifiIwasawaNotes}.

The second zero in (SP.1) is $X_\theta=0$, for the different
even character in (SP.3). Together with cyclicity of $X_\chi$,
it gives the integral Sharifi-map surjectivity by
\cite[Corollary~5.1.2 and its proof, p.~11]{WakeWangErickson2017}.
The primitive-character and exceptional-character assumptions
in its \S2.1 are satisfied because $e\ne0,2$. If $X_\chi=0$,
then $T_e=0$ and there is nothing further to prove. Otherwise
the cited theorem applies to this nonzero component.
Fukaya--Kato, \S\S7.1.8--7.1.10, turn that surjectivity into
generation of the finite-level $\theta$ target by cyclotomic
symbols, at every tower level
\cite[p.~86]{FukayaKato2012}.

We spell out why the characterwise version of
Proposition~\ref{prop:universal-pair} suffices. Its actual
class-two extension can be projected to the component $e$:
\[
0\longrightarrow
 (\widehat{\bigwedge}^{\,2}Q)^{[e]}
 \longrightarrow E^{[e]}
 \longrightarrow T_e\longrightarrow0.
\]
Its dual cohomological component has character $2-e$, exactly
$\theta$ in (SP.3). The norm-compatible cyclotomic inputs
annihilate all of $T$ by Lemma~\ref{lem:norm-torsion}, even
when the omitted odd torsion is nonzero. The same-level product
formula for nonunit exponents, parity reduction, and the
degree-one/degree-two comparison in
Lemma~\ref{lem:etale-prop} therefore supply the required
generating cups in this one target component. At $F_n$, use
Fukaya--Kato's coefficient index $n+1$ and then reduce to
coefficients $p^n$; this is the cofinal diagonal used in that
proposition.

If all quadratic companion coefficients vanished on a root
divisor for $T_e$, the split-socle construction in the proof of
Proposition~\ref{prop:universal-pair} would give a nonzero
$\Gamma$-equivariantly split socle in $T_e$. The explicit cup
primitive (NC.10) shows that one fixed power of $p$ kills the
image of every generating cup in its dual quotient, at every
level. Formula (NC.13), with the quotient transition maps,
would then make that same power kill a nonzero torsion-free
inverse-limit dual line. This is impossible. This contradiction
uses only the component $2-e$ of the target and only cyclicity
of $T_e$; no condition on another even component is used.
Thus Pair-NV holds for the actual free-$Q$ companion in every
selected component.

For $e=0,2$, the corresponding unramified odd components are
$X_\omega,X_{\omega^{-1}}$, both zero by the earlier
Stickelberger exception calculation. The Kummer
pseudoisomorphism gives $T_e[1/p]=0$, and $\mathbf Z_p$-torsion
freeness gives $T_e=0$.

\par\medskip\noindent\textbf{The actual subgroup and the linear coordinates.}
The compact splitting and eigenlift construction are the ones
proved in Theorem~\ref{thm:odd-torsion-cubic}, before any
cyclicity or nonvanishing assumptions are used. A retraction
that fixes the selected free generators and sends every omitted
generator to one is continuous on their pointed convergent
set. The universal property of the compact free basis extends
it to $H$. Its image is exactly $H_I$, a closed free factor.

Assign degrees one to all $Q$ variables, two to all even
torsion variables, and three to all odd torsion variables in
the full compact free basis. The weighted Magnus expansion is
injective. Every element of $H_I$ involves only the degree-one
variables and the selected degree-two variables. There is no
assumption that $\Gamma$ preserves $H_I$.

Let $q_-(z)$ be the characteristic polynomial of $\gamma$ on
$T^-[1/p]$, with $q_-=1$ if this module is zero. Remove from
the rational free target the finitely many generators $s_m$
for which $q_-(u^m)=0$, where $u=\chi_p(\gamma)$. Finiteness
follows because the powers $u^m$ are distinct. Write $\rho'$
for this quotient of $\rho$. Its target remains free on all
other odd generators, and its image remains Zariski dense in
every finite weight quotient. Every retained linear coordinate
kills $T^-$ by its annihilating polynomial and kills $T^+$
by parity. It therefore factors through $Q$.
The logarithm of an exact odd torsion eigenlift has no even
target lengths. After this linear-coordinate removal its
logarithm, and hence its group image, starts in length at least
three. Thus the full assigned filtration maps
to the target length filtration, including the omitted
generators that may occur in actual conjugates.

After a bounded rational change to an $R$-basis of $Q[1/p]$,
each such coordinate in its matching character is a nonzero
scalar multiple of evaluation at $T=u^m-1$. It is nonzero by
the stated density. Rescale the retained target generators to
make these singleton kernels exactly evaluations. This is an
automorphism of the weight-completed target, since each finite
weight quotient involves finitely many generators.

\par\medskip\noindent\textbf{Selected companion coordinates.}
For a nonzero selected component, choose an integral multiple
of a rational cyclic vector and an actual eigenlift $y_e$.
Let $y_{e,j}=\gamma^j(y_e)$ for
$0\leq j<\deg q_e$, where $q_e$ is its characteristic and
minimal polynomial. In assigned degree two, these elements
have an invertible rational linear combination of the old
$T_e$ variables, plus quadratic expressions in $Q$.
There is no linear contribution from another character
component. Replacing the selected degree-two variables by
these leading expressions is therefore an invertible
homogeneous triangular change, with bounded denominators in
each fixed degree. It preserves the algebra generated by $Q$
and the selected even variables. The actual elements $y_{e,j}$
need not belong to $H_I$; only their leading expressions are
used as this coordinate change.

The actual companion error has abelian class zero. Its degree-two
part is therefore a quadratic expression $C_e$ in the free
$Q$ variables alone. For $z_i=1+T_i$, comparison with the
target length-two coefficients gives
\[
 q_e(z_1z_2)B_{e,0}=C_e(T_1,T_2),\qquad
 B_{e,j}=(z_1z_2)^jB_{e,0}.
\]
The omitted generators have no role in this degree-two
identity. After a constant partial-fraction change, the pair
kernels have the form
\[
              \frac{C_{e;a,b}(T_1,T_2)}{(z_1z_2-\beta)^r},
 \qquad 1\leq r\leq\operatorname{mult}_{q_e}(\beta),
 \tag{SP.4}
\]
where $a+b=e$ in $\mathbf Z/(p-1)$ and wrong-character kernels
are zero. Pair-NV gives a nonzero restriction of at least one
matching numerator on each divisor $z_1z_2=\beta$.

\par\medskip\noindent\textbf{Injectivity on the selected free factor.}
Suppose $g\in H_I$ is nonidentity and $\rho'(g)=1$.
Take its first nonzero assigned Magnus degree $d$ after the
bounded changes above. At target length $d$, only this source
degree contributes: singleton images start in length one and
the even images in length two. Its coefficient is a finite
sum over ordered tilings of $d$ positions by singletons and
pairs, with the kernels (SP.4). The coefficient of a fixed
tiling is a bounded-denominator power series in its singleton
variables. Continuity and compactness justify these coefficient
evaluations by finite-basis truncation, exactly as in the
bounded substitution argument of
Theorem~\ref{thm:odd-torsion-cubic}.

Use sufficiently large odd test weights in each prescribed
residue character. These weights avoid the finitely many
discarded singleton resonances and every integer pair-sum
resonance of the selected $q_e$. They still have Cartesian
$p$-adic accumulation grids. Clearing all pair denominators
and applying the one-variable identity theorem successively
shows that the coefficient sum is identically zero in the
fraction field of the bounded power-series ring.

Among its nonzero tilings choose one with the largest number
of pairs, and then a componentwise maximal vector of pole
orders for fixed pair positions and root labels. On each
selected pair choose the odd residue types supplied by
Pair-NV. These choices are independent because the pairs
are disjoint. Set $S_B=z_i z_{i+1}-\beta_B$ on each selected
pair $B$. Retain its first variable, together with every
singleton variable, as base variables. Every denominator at
a different adjacent pair has nonzero specialization in the
fraction field of these base variables: an interval of length
two can be a union of selected disjoint pairs only by being
one of them. A different root at the same pair also gives a
unit. Thus the identity has a multivariate Laurent expansion
with finite negative bounds in the selected $S_B$.

Extract the coefficient of the selected product of negative
powers. A contributing term must contain every selected pair
and root. Maximality of the number of pairs forbids extras;
maximality of the pole-order vector excludes every strict
componentwise domination, and an incomparable vector misses
one required exponent. The selected odd residue types identify
the even character, and cyclicity excludes several distinct
blocks with the same character and root. Hence the only
possible contribution is the selected nonzero singleton
coefficient multiplied by the nonzero restrictions of the
selected $C_e$. These are functions in independent base
variables, so their product is nonzero in a field. This
contradicts the zero Laurent coefficient. With no pairs the
same argument is just the analytic identity theorem.

It follows that $\rho'|_{H_I}$ is injective. Since $\rho'$
is a quotient of $\rho$, (SP.2) follows. In particular the
original outer action is faithful on $H_I$, because an
element acting trivially there also has trivial rational image.
\end{proof}

\begin{corollary}[K17-SELECTED-ARITHMETIC-RANGE]
\label{cor:selected-arithmetic-range}
Under the stipulated arithmetic theorem, for every $p\geq5$
one may take in Theorem~\ref{thm:selected-pair-faithfulness}
\[
 I=\{[-2j],[2j+2]:0\leq j\leq12\}
        \ \subset\ \mathbf Z/(p-1).
 \tag{SP.5}
\]
Thus the actual quadratic companions in these even characters
are rationally cyclic and satisfy Pair-NV, and the closed free
factor generated by all free odd lifts and these even torsion
lifts is faithful. The companion conditions in the omitted
characters, including actual cubic nonvanishing, remain separate.
\end{corollary}

\begin{proof}
For $e=-2j$, Corollary~\ref{cor:twelve-degrees-period} and
Proposition~\ref{prop:componentwise-iwasawa-control} give
$X^{[e]}=0$, using the trivial component when $j=0$, while
Proposition~\ref{prop:positive-iwasawa-components} gives
$X^{[2-e]}=X^{[2j+2]}=0$. The involution $e\mapsto2-e$
interchanges the two classes displayed for each $j$, so the
same two facts verify (SP.1) for $e=2j+2$. This includes all
coincidences modulo $p-1$ and the empty exceptional companions.
Apply the theorem. More generally the same conclusion holds
with $12$ replaced by any $R\leq15$ once the vanishings
$K_{4j}(\mathbf Z)=0$ have been established for $1\leq j\leq R$,
because $2R+2\leq32$ lies within the positive-character range.
\end{proof}

This restriction theorem does not identify the kernel on all of
$H$. The omitted even and odd torsion generators and words
involving them are outside its assertion. The actual subgroup
is a specified free factor rather than a canonical
$\Gamma$-submodule of the Galois group.
